\documentclass[12pt,a4paper]{article}

\begin{document}

\section*{Chapter 1: Introduction}
Drug discovery is a complex and iterative process that benefits significantly from computational methods capable of rapidly evaluating large chemical spaces. This chapter introduces the role of molecular modelling in supporting early-stage decision-making, highlighting the complementary strengths of human reasoning and computational efficiency. While computers excel at executing repetitive, physics-based calculations, human expertise remains essential for interpreting results and guiding scientific strategy. The chapter outlines how molecular docking and related techniques rely on simplifying assumptions, such as approximate treatment of solvation and entropy, to achieve tractable computation, while also emphasising the limitations introduced by these approximations. Fundamental concepts of molecular recognition are discussed, including the role of non-covalent interactions such as hydrogen bonding, electrostatics, van der Waals forces, and hydrophobic effects in determining protein–ligand binding. Binding affinity is introduced as a measurable quantity linked to thermodynamic parameters and experimental observables such as dissociation constants. The chapter further highlights the practical utility of computational screening in reducing experimental costs by prioritising promising candidates. Finally, it positions the book as a practical guide designed to enable readers, particularly beginners, to apply molecular docking methods effectively while maintaining a clear understanding of their theoretical basis, assumptions, and limitations.

\section*{Chapter 2: Conformational Search Methods for Molecular Docking}
This chapter presents the principles and computational strategies used to explore molecular conformational space in the context of molecular docking. It begins by establishing the importance of molecular conformation in determining chemical and biological properties, particularly binding affinity. The combinatorial complexity associated with flexible molecules is discussed, demonstrating how the number of possible conformations increases exponentially with the number of rotatable bonds. To address this challenge, the concept of the potential energy surface is introduced as a mathematical representation of molecular energy as a function of atomic coordinates. The chapter explains how local minima correspond to stable conformations and how energy barriers influence conformational transitions. The use of internal coordinates, including bond lengths, angles, and torsions, is described as a practical framework for molecular mechanics modelling. Key mathematical tools, such as gradients and the Hessian matrix, are introduced to classify stationary points and assess stability. Simplified representations, such as one-dimensional energy profiles along reaction coordinates, are also discussed to aid interpretation. Overall, the chapter provides a conceptual and practical foundation for understanding how docking algorithms efficiently identify bioactive conformations while balancing computational cost and accuracy.

\section*{Chapter 3: Scoring Functions}
This chapter develops the thermodynamic framework underlying scoring functions used in molecular docking. Protein–ligand binding is described as a reversible equilibrium process characterised by association and dissociation constants, which are directly related to the standard free energy of binding. The relationship between free energy, enthalpy, and entropy is presented to provide a physical basis for interpreting molecular interactions. The chapter highlights the sensitivity of binding affinity to small changes in free energy, illustrating how modest energetic differences can translate into large variations in experimentally observed constants. Experimental approaches for measuring thermodynamic parameters are discussed, alongside the challenges faced in computational estimation. Particular emphasis is placed on the interplay between enthalpic and entropic contributions, including the phenomenon of enthalpy–entropy compensation, which complicates intuitive interpretation. The chapter also introduces covalent binding as a two-step process involving initial non-covalent association followed by chemical bond formation. By linking thermodynamic principles to computational approximations, the chapter provides a rigorous foundation for understanding how scoring functions operate and why their predictions must be interpreted with appropriate caution.

\section*{Chapter 4: Handling Biomacromolecular Structures for Molecular Docking}
This chapter examines the role of biomacromolecular structure quality in determining the reliability of molecular docking results. It emphasises that experimentally derived structures are models that require careful validation and interpretation. The chapter introduces key structural determination techniques, including X-ray crystallography, cryo-electron microscopy, and nuclear magnetic resonance, and discusses their respective strengths and limitations. Important quality metrics, such as resolution and agreement factors, are described in relation to their impact on structural accuracy, particularly within binding sites. Common issues, including missing residues, incorrect oligomeric states, poorly modelled ligands, and crystallisation artefacts, are highlighted as potential sources of error in docking workflows. The chapter also considers the role of computationally predicted structures and their associated uncertainties. Practical guidance is provided on preparing protein structures, including identifying biologically relevant assemblies, validating ligand geometry, and assessing structural flexibility. Overall, the chapter underscores that careful selection, validation, and preparation of biomacromolecular structures are essential prerequisites for obtaining meaningful and reliable docking results.

\section*{Chapter 5: Methods to Identify Binding Sites on Proteins}
This chapter explores computational approaches for identifying ligand-binding sites on proteins, a critical step preceding molecular docking. Binding sites are defined in both structural and energetic terms, with emphasis on their role in governing molecular recognition and specificity. Different classes of binding sites, including orthosteric, allosteric, and transient or cryptic pockets, are introduced to illustrate the diversity of interaction environments. The chapter reviews a range of computational methods, from geometry-based approaches that detect cavities on protein surfaces to machine learning techniques that integrate sequence, structural, and dynamic information. Key descriptors, such as pocket size, shape, hydrophobicity, and polarity, are discussed in the context of characterising binding regions. The concept of druggability is introduced as a quantitative measure of a site’s suitability for small-molecule binding. Thermodynamic considerations are also presented to link binding site properties with free energy contributions. The chapter emphasises validation strategies to distinguish biologically relevant sites from artefacts and provides guidance for identifying challenging targets, including cryptic and transient pockets.

\section*{Chapter 6: Preparation of Small-Molecule Libraries}
This chapter outlines the theoretical and practical considerations involved in preparing small-molecule libraries for molecular docking. It begins by distinguishing between ligands and drugs and emphasising the importance of accurate molecular representation. The transformation of molecules from one-dimensional or two-dimensional formats into three-dimensional geometries suitable for docking is described. Physicochemical filtering criteria, including molecular weight, lipophilicity, and hydrogen bonding capacity, are introduced as methods for enriching libraries with drug-like compounds. The chapter details key preprocessing steps, including protonation state prediction, tautomer enumeration, stereochemical assignment, and conformational sampling. The importance of assigning appropriate force field parameters and partial charges is also discussed to ensure compatibility with docking algorithms. Quality control procedures are emphasised to eliminate chemically inconsistent or unrealistic structures that could compromise results. By integrating theoretical principles with practical workflows, the chapter demonstrates how careful ligand preparation significantly enhances the reliability and efficiency of virtual screening campaigns.

\section*{Chapter 7: Molecular Docking Methods and Programs}
This chapter provides an overview of molecular docking methodologies and commonly used software tools. It introduces the concept of bioactive conformations and explains their central role in predicting ligand binding. The two fundamental components of docking workflows—conformational sampling and scoring—are described in detail. Different algorithmic strategies, including geometric matching, incremental construction, and stochastic search methods, are discussed to illustrate how docking programs explore ligand orientations and conformations within a binding site. The role of scoring functions in ranking poses based on interaction energies is also examined, along with their inherent limitations in predicting absolute binding affinity. Practical aspects of docking workflows are covered, including receptor preparation, binding site definition, ligand preparation, parameter selection, and validation. The chapter emphasises the importance of understanding software-specific parameters and consulting validation studies to ensure appropriate application. Overall, it provides a structured framework for selecting and using docking programs effectively in drug discovery.

\section*{Chapter 8: Analysing Results from Molecular Docking}
This chapter addresses the interpretation and validation of molecular docking results, highlighting analysis as a critical step in the workflow. Docking outputs are described as comprising predicted binding poses and associated scores, which serve as approximations rather than direct measures of binding affinity. The distinction between pose prediction accuracy and ranking reliability is emphasised. The chapter provides practical guidance on evaluating results through visual inspection, assessment of key interactions, and identification of unrealistic features such as steric clashes or strained conformations. Additional outputs, including interaction fingerprints and pose clustering, are discussed as tools for deeper analysis. Common pitfalls and misleading scenarios are presented to help users avoid overinterpretation. The chapter reinforces the view that docking should be treated as a hypothesis-generating method and that results must be validated through complementary computational or experimental approaches.

\section*{Chapter 9: Molecular Dynamics for Post-processing Docking Results}
This chapter introduces molecular dynamics simulations as a complementary approach to refine and validate docking predictions. It outlines the key limitations of docking, including simplified treatment of receptor flexibility and approximate scoring functions. Molecular dynamics is presented as a method that captures the time-dependent behaviour of protein–ligand systems in explicit solvent environments, providing a more realistic representation of interactions. The theoretical framework based on classical mechanics is described, including equations of motion and statistical ensembles. Practical aspects of simulation, such as numerical integration and the use of thermostats and barostats, are also discussed. The chapter highlights how molecular dynamics can be used to assess the stability of docking poses, explore conformational changes, and improve binding predictions. Enhanced sampling techniques are introduced as methods to overcome limitations in sampling rare events. Overall, the chapter demonstrates how molecular dynamics enhances the interpretability and reliability of docking results.

\section*{Chapter 10: Estimating Binding Affinity by Free-Energy Methods}
This chapter presents advanced computational approaches for estimating protein–ligand binding affinity using statistical mechanics. It begins by highlighting the limitations of conventional docking scoring functions and motivates the use of more rigorous free energy methods. The concept of the partition function is introduced to provide a statistical description of molecular systems and their thermodynamic properties. The chapter describes various classes of methods, including end-state approaches, linear interaction energy models, and alchemical techniques such as free energy perturbation and thermodynamic integration. The use of thermodynamic cycles to compute binding free energies indirectly is explained. The distinction between relative and absolute binding free energy calculations is also discussed in the context of drug discovery applications. While these methods offer improved accuracy, their computational cost and complexity are acknowledged. The chapter positions free energy calculations as complementary tools to docking for refining predictions and supporting lead optimisation.

\section*{Chapter 11: Communicating Results from Molecular Docking}
This chapter provides guidance on presenting molecular docking results effectively to both expert and non-expert audiences. It emphasises the importance of clarity, context, and appropriate interpretation when communicating computational findings. The distinction between binding poses and docking scores is explained, along with their respective meanings and limitations. The chapter highlights the need to compare results with known reference compounds to provide meaningful context. It also addresses common misconceptions, such as overinterpreting small differences in docking scores. The importance of reporting methodological details, including software, parameters, and validation procedures, is emphasised to ensure reproducibility and credibility. Practical recommendations are provided for structuring manuscripts and presentations. The chapter concludes by reinforcing that docking results should be viewed as hypotheses that guide experimental validation rather than definitive conclusions.

\section*{Chapter 12: FAIR and Reproducible Reporting of Molecular Docking Studies}
This chapter discusses the importance of reproducibility and transparency in computational docking studies, focusing on the application of FAIR principles. It begins by outlining challenges associated with reproducibility, including limited access to data and insufficient methodological reporting. The FAIR principles—Findable, Accessible, Interoperable, and Reusable—are introduced as a framework for improving research practices. The chapter explains how these principles can be applied to docking workflows through proper documentation of datasets, software, parameters, and protocols. The role of persistent identifiers, metadata, and standardised formats is emphasised in enhancing discoverability and usability. Practical considerations, such as data repositories and access conditions, are also discussed. By promoting transparent and well-documented workflows, the chapter highlights how adherence to FAIR principles improves scientific credibility and facilitates collaboration.

\section*{Chapter 13: Union of Artificial Intelligence and Molecular Docking}
This chapter examines the integration of artificial intelligence (AI) with molecular docking, highlighting how data-driven approaches are transforming traditional computational workflows in drug discovery. Classical docking methods, based on heuristic conformational sampling and approximate scoring functions, are first revisited to outline their limitations, including simplified energy models, restricted treatment of protein flexibility, and challenges in exploring high-dimensional conformational spaces. The chapter then introduces the foundations of deep learning, including neural network architectures such as convolutional and graph-based models, which enable the learning of complex, non-linear relationships from structural data. AI-based scoring functions are discussed as flexible alternatives to empirical models, capable of improving pose ranking and binding affinity prediction. The emergence of geometric deep learning approaches that directly predict ligand binding poses, bypassing exhaustive sampling, is also described. Diffusion-based generative models and equivariant neural networks are presented as key advances that enhance both efficiency and accuracy in docking. The chapter further explores the impact of breakthroughs in protein structure prediction, particularly AI-based models, which provide high-quality structural inputs for docking workflows. More advanced co-folding models that simultaneously predict protein–ligand complexes and binding affinity are introduced as a unifying framework that bridges structure prediction and docking. Despite these advances, the chapter emphasises challenges related to data quality, model generalisation, and the need for physics-based validation. Overall, it presents AI as a powerful complement to traditional methods, with the potential to significantly improve the accuracy, speed, and scope of molecular docking in modern drug discovery.

\section*{Chapter 14: Special Topics in Molecular Docking}
This chapter explores advanced applications of molecular docking that extend beyond conventional single protein–ligand systems. It focuses on complex scenarios such as covalent docking, multi-ligand binding, and the modelling of ternary complexes. Covalent docking is described as a process that incorporates chemical reactions into binding predictions, requiring specialised algorithms and scoring functions. The chapter also discusses approaches for modelling cooperative or competitive binding of multiple ligands within a single binding site. In the context of emerging therapeutic strategies, methodologies for designing bifunctional (PROTACS) molecules that induce protein–protein interactions are introduced. Various computational approaches, including constraint-based and hybrid quantum-mechanical methods, are described to address these challenges. The chapter highlights both the opportunities and limitations associated with these advanced applications, emphasising the need for careful validation and interpretation.

\end{document}